\chapter{Enquadramento}

Este capítulo está organizado em duas secções. A primeira descreve o contexto 
da gestão de clubes desportivos de menor dimensão e a realidade específica do 
clube para o qual a aplicação foi desenvolvida. A segunda apresenta soluções 
existentes neste domínio e posiciona a aplicação JagozApp face a essas 
soluções.

\section{Gestão de Clubes Desportivos de Pequena Dimensão}

A gestão de um clube desportivo de pequena ou média dimensão envolve um 
conjunto alargado de tarefas administrativas, financeiras e de comunicação que, 
com frequência, são asseguradas por um número reduzido de pessoas, muitas vezes 
em regime de voluntariado. Processos como a gestão de sócios, a inscrição de 
atletas, o controlo de pagamentos ou a angariação de patrocínios tendem a estar 
dispersos por folhas de cálculo, formulários em papel, correio eletrónico e 
comunicações informais.

No caso concreto do Grupo Desportivo União Ericeirense, identificou-se uma 
necessidade que ultrapassa a gestão interna: a aproximação entre o clube e a 
sua comunidade. A ausência de meios digitais orientados para o exterior 
traduzia-se em oportunidades perdidas, tanto na angariação de novos sócios como 
de patrocinadores, sobretudo junto de pessoas com pouca disponibilidade para se 
deslocarem presencialmente à sede do clube.

\section{Sistemas Semelhantes}

Existem diferentes tipos de soluções relacionadas com a presença digital de um 
clube desportivo, que podem ser agrupadas em duas categorias.

A primeira categoria corresponde às \textbf{plataformas de gestão de clubes}, 
das quais o EMJOGO~\cite{emjogo} é um exemplo representativo no contexto 
português. Trata-se de uma solução comercial, disponibilizada no modelo de 
\textit{software} como serviço (SaaS), que centraliza numa única plataforma a 
gestão desportiva, administrativa e financeira do clube, incluindo a gestão de 
atletas, sócios, mensalidades e patrocinadores. O EMJOGO disponibiliza ainda um 
\textit{website} e uma aplicação de comunicação com a comunidade do clube.

A segunda categoria corresponde à \textbf{presença \textit{online} própria de 
cada clube}, que, no caso dos clubes distritais de menor dimensão, se traduz 
maioritariamente em \textit{websites} de carácter institucional e informativo. 
Estes apresentam a história do clube, as suas equipas e notícias, mas raramente 
disponibilizam funcionalidades interativas que permitam ao adepto realizar 
ações diretamente, como aderir como sócio, inscrever um atleta ou adquirir 
bilhetes.

É neste contexto que o clube apresentou a sua proposta inicial. O Grupo 
Desportivo União Ericeirense já utiliza o EMJOGO para a sua gestão interna, e o 
objetivo inicialmente colocado era o de estabelecer uma ligação entre a nova 
aplicação e essa plataforma, de forma a evitar a duplicação manual de dados. 
Contudo, o EMJOGO é uma solução comercial de código fechado e não disponibiliza 
publicamente uma interface de integração que o permitisse. Por esse motivo, uma 
integração direta não se revelou viável no âmbito deste projeto.

Perante esta limitação, o foco do projeto foi redefinido. Em vez de procurar 
centralizar ou integrar a gestão já assegurada pelo EMJOGO, a aplicação 
JagozApp passou a ter como objetivo complementar essa plataforma, oferecendo ao 
clube aquilo que considera ser a sua principal lacuna no contexto específico de 
um clube de pequena dimensão: uma ligação direta e acessível entre o clube e o 
adepto. Assim, a JagozApp distingue-se por reunir, numa aplicação 
\textit{web} desenvolvida à medida do clube, um conjunto de funcionalidades 
voltadas para o exterior — a adesão de sócios, a inscrição de atletas, o registo 
de patrocinadores e a venda de bilhetes com validação por código QR. Ao 
contrário de uma solução comercial genérica, esta é uma aplicação própria do 
clube, sobre a qual este detém controlo total e que pode evoluir de acordo com 
as suas necessidades futuras.